\appendix
\section{Scope, Edition, and Qualifications}

% This is adult-facing audit matter, not a learner lesson.  A compact local
% two-column setting keeps the qualifications, bibliography, generation
% metadata, and rights colophon together on the terminal page without changing
% the type or spacing used anywhere else in the guide.
\begingroup
\footnotesize
\microtypesetup{protrusion=false,expansion=false}
\RaggedRight
\setlength{\columnsep}{0.8em}
\setlength{\parindent}{0pt}
\setlength{\parskip}{0.22em}
\newcommand{\TerminalNote}[2]{\textbf{#1.} #2\par}
\begin{multicols}{2}

\TerminalNote{Form}{\GuideFormStatement}

\TerminalNote{Training boundary}{\GuideBoundaryStatement}

\TerminalNote{Roster and model}{\GuideRosterStatement}

\TerminalNote{Text and pronunciation}{The Latin response sequence is collated to
the 1962 Vatican typical Missal. The printed house form expands the ligatures
\textit{ae} and \textit{oe}; stress marks and syllable divisions appear only
on separately labeled learner lines. The pronunciation course teaches the
Roman-style ecclesiastical Latin described in the 1961 \textit{Liber Usualis}.
Its broad IPA, sound spelling, memory aids, and short English meanings are
original editorial layers. They are not chant notation or an English ritual
translation. For sung delivery, the appointed chant and the choir director
control melody, rhythm, and phrasing.}

\TerminalNote{Edition limits}{This ordinary route excludes the Asperges,
processions, Candlemas, Ash Wednesday, Palm Sunday, Holy Week, Rogation rites,
Nuptial and Requiem ceremonies, funerals and absolutions, Mass before the
Blessed Sacrament exposed, pontifical functions, Benediction, and Communion
outside Mass. Compact branch warnings do not teach those functions. The
Leonine prayers are outside Mass and are not included. The printed 1962
Communion rite has no second \textit{Confiteor}.}

\columnbreak
\TerminalNote{Authority and training}{This is a study and rehearsal aid, not an
official liturgical book, parish directive, certificate of competence, or
statement of present authorization. The priest and master of ceremonies govern
the actual function. \GuideSafetyStatement}

\TerminalNote{Review and rights}{Internal source, text, layout, and production
checks are recorded in the series research files. Independent liturgical,
ceremonial, Latin-pronunciation, pedagogical, rights, and ecclesiastical review
remain outstanding. Project-created teaching matter falls within the
repository license boundary; received liturgical text, official sources,
chant, scans, fonts, and third-party manuals retain their own status.
Installation does not authorize public distribution.}

% Missa Cantata already has the series' longest learner-facing contents map.
% Keep its technical bibliography and metadata in the book without making a
% nearly empty second contents page for those two entries alone.
\ifMissaCantataGuide
  \addtocontents{toc}{\protect\setcounter{tocdepth}{0}}
\fi
\titleformat{\section}{\normalsize\bfseries}{\thesection}{0.45em}{}
\titlespacing*{\section}{0pt}{0.45\baselineskip}{0.2\baselineskip}
\section{References}

\begin{itemize}[leftmargin=1.05em,labelsep=0.35em,topsep=0.1em,
  itemsep=0.12em,parsep=0pt,partopsep=0pt]
  \item \textit{Missale Romanum}, editio typica (Typis Polyglottis Vaticanis,
  1962), especially the General Rubrics, \textit{Ritus servandus}, and
  \textit{Ordo Missae}; page-image facsimile:
  \url{https://media.churchmusicassociation.org/pdf/missale62.pdf}.
  \item Sacred Congregation of Rites, \textit{De musica sacra et sacra
  liturgia}, 3 September 1958, \textit{AAS} 50 (1958), 630--663, especially
  nn.\ 24--27 and 31--32:
  \url{https://www.vatican.va/archive/aas/documents/AAS-50-1958-ocr.pdf}.
  \item \textit{Liber Usualis} (Solesmes, 1961), English introduction,
  section IX, pp.\ xxxv--xxxix.
  \item \GuideCeremonialReferences
\end{itemize}

Full source roles, access state, discrepancies, and rights treatment are
recorded in the series edition manifest and source audit.

\end{multicols}
\endgroup
